\documentclass{article}

\usepackage{amsmath, amsthm, amssymb, amsfonts}
\usepackage{thmtools}
\usepackage{graphicx}
\usepackage{setspace}
\usepackage{geometry}
\usepackage{float}
\usepackage{hyperref}
\usepackage[utf8]{inputenc}
\usepackage[english]{babel}
\usepackage{framed}
\usepackage[dvipsnames]{xcolor}
\usepackage{tcolorbox}

\colorlet{LightGray}{White!90!Periwinkle}
\colorlet{LightOrange}{Orange!15}
\colorlet{LightGreen}{Green!15}

\newcommand{\HRule}[1]{\rule{\linewidth}{#1}}

\declaretheoremstyle[name=Theorem,]{thmsty}
\declaretheorem[style=thmsty,numberwithin=section]{theorem}
\tcolorboxenvironment{theorem}{colback=LightGray}

\declaretheoremstyle[name=Proposition,]{prosty}
\declaretheorem[style=prosty,numberlike=theorem]{proposition}
\tcolorboxenvironment{proposition}{colback=LightOrange}

\declaretheoremstyle[name=Principle,]{prcpsty}
\declaretheorem[style=prcpsty,numberlike=theorem]{principle}
\tcolorboxenvironment{principle}{colback=LightGreen}

\setstretch{1.2}
\geometry{
    textheight=9in,
    textwidth=5.5in,
    top=1in,
    headheight=12pt,
    headsep=25pt,
    footskip=30pt
}

% ------------------------------------------------------------------------------

\begin{document}

% ------------------------------------------------------------------------------
% Cover Page and ToC
% ------------------------------------------------------------------------------

\title{ \normalsize \textsc{}
		\\ [2.0cm]
		\HRule{1.5pt} \\
		\LARGE \textbf{\uppercase{Relatório do trabalho\\ Minimundo Livraria}
		\HRule{2.0pt} \\ [0.6cm] \LARGE{Ferramentas da Internet} \vspace*{10\baselineskip}}
		}
\date{}
\author{\textbf{Autores:} \\
		Clara Cotta\\
		Deborah Martins\\
        Manuela Brumatti\\
        Yasmin Silveira\\
        Ysadora Silva\\
		23/04/2026}

\maketitle
\newpage

% ------------------------------------------------------------------------------

\section{Introdução}
Este minimundo descreve o funcionamento de uma livraria, envolvendo o controle de livros, clientes, vendas e fornecedores. O sistema concede cadastrar livros contendo informações, registrar clientes e realizar compras. Também garante o controle de estoque e o acompanhamento das vendas, ajudando a manter tudo atualizado. Dessa forma, possui o objetivo de armazenar, organizar e gerenciar de forma eficiente as as principais atividades da livraria, contribuindo para um melhor funcionamento do sistema.
Sua importância está voltada em proporcionar maior controle e confiabilidade nas operações da livraria, contribuindo para a redução de instabilidades, melhor organização das informações e agilidade nos processos internos. Nesse contexto, a ordenação abrange o acompanhamento das principais atividades da livraria, incluindo o relacionamento entre os dados envolvidos nas vendas, no cadastro de pessoas, no fornecimento de produtos e na logística de atendimento.

\section{Minimundo}
A livraria Books Space precisa de um sistema de banco de dados integrado para gerenciar seu estoque, vendas, equipe de funcionários e a logística com fornecedores, controlando todo o fluxo desde a entrada do livro até a venda ao consumidor final. No centro do sistema, utiliza-se uma estrutura de especialização onde a entidade Pessoa, composta por CPF (ID), nome, e-mail e endereço, ramifica-se em Funcionário, que possui matrícula, cargo, salário e data de admissão; e Cliente, que armazena data de nascimento, pontos de fidelidade e data do último login. Um cliente pode realizar vários pedidos, enquanto um funcionário pode atender diversos clientes, mas cada pedido pertence a apenas um cliente e um funcionário. O acervo é representado pela entidade Livro, que possui ISBN (ID), título, preço de venda e ano de publicação, estando associado ao controle de Estoque, responsável por controlar o código de armazenamento (ID), registrar a quantidade de livros e o espaço disponível. Cada livro está vinculado, obrigatoriamente, a uma única Editora, que possui CNPJ (ID), nome e contato, e a uma única Categoria, que contém identificação (ID), nome e descrição. Já a relação entre Livro e Autor é de muitos para muitos, permitindo que uma obra tenha múltiplos escritores e que um autor possa estar associado a diversas obras. O processo de venda é representado pela entidade Pedido de Venda, que possui número identificador (ID), data, valor total e status. Cada pedido está associado a um único cliente e a um único funcionário responsável pelo atendimento. Um pedido de venda pode conter vários livros, sendo preciso registrar a quantidade de cada item adicionado. Assim, estabelece-se uma relação entre Pedido de Venda e Livro, em que um pedido pode conter muitos livros e um mesmo livro pode estar presente em vários pedidos. A entidade Fornecedor é responsável por gerenciar as informações dos distribuidores de livros da livraria. Ela é composta por CNPJ (ID), nome, contato e endereço. Um fornecedor pode fornecer diversos livros e um mesmo livro pode ser adquirido de diferentes fornecedores. A livraria registra a entrada de livros por meio da entidade Nota de Entrada, formada por número (ID), data e valor total. Cada nota pode conter vários Itens de Entrada, sendo que cada item está ligado a um único Livro, registrando quantidade e custo por unidade. A entidade Livro pode estar associada a diferentes fornecedores, permitindo que um mesmo livro seja adquirido de mais de uma origem. O registro da entrada atualiza o Estoque e pode estar associado a um Funcionário, que fica responsável pelo recebimento. Além disso, um Fornecedor pode estar vinculado a diversas notas de entrada. 
 
\section{Divisão de tarefas}
Do total de 10 entidades presentes no minimundo, Clara, Leticya e Ysadora fizeram as 7 primeiras, e o restante foi dividido entre Deborah, Manuela e Yasmin. Clara e Deborah ficaram responsáveis pelo relatório, montando a introdução.



\newpage


\end{document}
